\section{An anchor for successive choices}
\label{sec:anchor}

We construct a point $x\in K$ and a family of nonexpansive maps whose
images of $x$ have two useful properties.  Their fixed point defects can
be made arbitrarily small, even after additional requirements have been
imposed by earlier choices, and finite convex combinations of these
images can return arbitrarily close to $x$ in norm.  The construction
uses compactness in a space of maps; finite simultaneous approximation
in norm will justify the compositions that occur there.

\subsection{Resolvents and displacement estimates}

We begin with the classical resolvent construction; see
\citet[pp.~75--76]{Goebel1975}.

\begin{lemma}[Resolvents]
\label{lem:resolvents}
For every $n\geq 1$ and $a\in K$, there is a unique point $R_n(a)\in K$
satisfying
\begin{equation}
\label{eq:resolvent-definition}
 R_n(a)=\frac{a}{n}+\left(1-\frac1n\right)T R_n(a).
\end{equation}
The map $R_n:K\to K$ is nonexpansive, and
\begin{equation}
\label{eq:resolvent-estimates}
 \norm{T R_n(a)-R_n(a)}\leq\frac1n,
 \qquad
 \norm{R_n(a)-a}\leq(n-1)\norm{Ta-a}
 \quad(a\in K).
\end{equation}
\end{lemma}

\begin{proof}
For $n=1$, the defining equation gives $R_1(a)=a$.  For $n>1$, the map
\[
 u\longmapsto \frac{a}{n}+\left(1-\frac1n\right)Tu
\]
takes the complete set $K$ into itself and has Lipschitz constant
$q=1-1/n<1$.  Starting from any point of $K$, its iterates have successive
displacements bounded by $q^r$ times the first displacement.  The
geometric series shows that the iterates are Cauchy.  Their limit solves
\eqref{eq:resolvent-definition} by norm continuity.  Two solutions have
distance at most $q$ times their distance and therefore coincide.

For $u=R_n(a)$ and $v=R_n(b)$, comparison of their equations gives
\[
 \norm{u-v}\leq\frac1n\norm{a-b}
              +\left(1-\frac1n\right)\norm{u-v},
\]
which proves nonexpansiveness.  The identity
$Tu-u=(Tu-a)/n$ and $\diam K=1$ give the first estimate in
\eqref{eq:resolvent-estimates}.  For the second, the defining equation
and nonexpansiveness of $T$ give
\[
 \norm{u-a}
 \leq\left(1-\frac1n\right)
       \bigl(\norm{u-a}+\norm{Ta-a}\bigr).
\]
Rearranging proves the claim.
\end{proof}

Let $\mathcal G$ be the family generated by the identity map and the
maps $R_n$ using finitely many compositions and pointwise finite convex
combinations.  Every member of $\mathcal G$ is a nonexpansive selfmap of
$K$.  Moreover, for each $G\in\mathcal G$ there is a finite constant
$c_G\geq0$ such that
\begin{equation}
\label{eq:map-displacement}
 \norm{G(a)-a}\leq c_G\norm{Ta-a}
 \qquad(a\in K).
\end{equation}
For the generators, take $c_{\mathrm{Id}}=0$ and $c_{R_n}=n-1$.
Convex combinations preserve the estimate with the same convex
combination of the constants.  For compositions, nonexpansiveness gives
\begin{align*}
 \norm{G_1(G_2(a))-a}
 &\leq\norm{G_1(G_2(a))-G_1(a)}+\norm{G_1(a)-a}\\
 &\leq\norm{G_2(a)-a}+\norm{G_1(a)-a}\\
 &\leq(c_{G_2}+c_{G_1})\norm{Ta-a}.
\end{align*}
This proves \eqref{eq:map-displacement} throughout $\mathcal G$.
The same composition estimate underlies Bruck's semigroup of maps
asymptotically fixing a common approximate fixed point sequence
\cite[p.~115]{Bruck1975}.

For $m\geq1$, define the tail family
\begin{equation}
\label{eq:tail-family}
 \mathcal H_m
   =\{R_n\circ G:n\geq m,\ G\in\mathcal G\}.
\end{equation}
These families are nonempty, contained in $\mathcal G$, and decreasing
with $m$.  Their outer resolvents give the uniform bound
\begin{equation}
\label{eq:tail-defect}
 \norm{T h(a)-h(a)}\leq\frac1m
 \qquad(h\in\mathcal H_m,\ a\in K).
\end{equation}

\subsection{A compact semigroup of maps}


Convex semigroups of nonexpansive maps, with the topology of pointwise
weak convergence, were studied by Bruck \cite{Bruck1975}. His
common-fixed-point theorem assumes strict convexity and pairwise
intersection of the closed convex hulls of the ranges of the maps.
Here the semigroup will supply an idempotent anchor, through Ellis's
lemma; the tree and switching arguments will then produce the fixed
point. The composition estimates below keep the needed continuity
properties explicit.

Give $K^K$, the set of all maps from $K$ into itself, the topology of
pointwise weak convergence.  It is a compact Hausdorff space by the
product compactness in Lemma~\ref{lem:weak-facts}.  We write
$\overline{\mathcal B}^{\,\mathrm{pw}}$ for closure in this space.

\begin{lemma}[Finite simultaneous approximation]
\label{lem:finite-tuples}
Let $\mathcal B\subseteq K^K$ be a nonempty convex family of
nonexpansive maps.  Every map
$h\in\overline{\mathcal B}^{\,\mathrm{pw}}$ is nonexpansive.
For any integer $q\geq1$, points $a_1,\ldots,a_q\in K$, and
$\delta>0$, there is a single
$g\in\mathcal B$ such that
\[
 \max_{1\leq i\leq q}\norm{g(a_i)-h(a_i)}<\delta.
\]
\end{lemma}

\begin{proof}
For each fixed pair $a,b\in K$, the condition
$\norm{h(a)-h(b)}\leq\norm{a-b}$ is closed in the pointwise weak
topology: evaluation at $a,b$ and subtraction are weakly continuous,
and norm balls are weakly closed.  Intersecting these conditions over
all pairs proves the first assertion.

For the second, consider the convex subset
\[
 V=\{(g(a_1),\ldots,g(a_q)):g\in\mathcal B\}\subseteq X^q.
\]
The tuple $(h(a_1),\ldots,h(a_q))$ belongs to its weak closure, by
continuity of evaluation.  The weak topology of the finite Banach
product $X^q$ is the product of its coordinate weak topologies: each
continuous linear functional on $X^q$ is the sum of its coordinate
restrictions.  By separation, the weak and norm closures of a convex
subset of $X^q$ agree.  Applying this fact with the maximum coordinate
norm gives one element of $V$ satisfying all the desired inequalities.
\end{proof}

We will also use the following elementary compact semigroup argument,
the right-translation version of \citet[Lemma~1, p.~402]{Ellis1958}.
Ellis uses the opposite convention for composition.
Here a semigroup is a set with an associative multiplication.

\begin{lemma}[Idempotent in a compact semigroup]
\label{lem:idempotent}
Let $S$ be a nonempty compact Hausdorff semigroup.  Suppose that, for
each $p\in S$, the right translation $h\mapsto hp$ is continuous.
Then $S$ contains an idempotent: an element $p$ with $p^2=p$.
\end{lemma}

\begin{proof}
The intersection of a chain of nonempty compact subsemigroups of $S$
is nonempty by compactness and is again a compact subsemigroup.
Zorn's lemma therefore gives a minimal nonempty compact subsemigroup
$S'\subseteq S$.  Fix $p\in S'$.  The set $S'p$ is a nonempty compact
subset of $S'$ by continuity of right translation.  It is a
subsemigroup because
\[
 (ap)(bp)=(apb)p\in S'p\qquad(a,b\in S').
\]
Minimality gives $S'p=S'$, so in particular there is an $h\in S'$ with
$hp=p$.  The set
\[
 E=\{h\in S':hp=p\}
\]
is thus nonempty.  It is compact by continuity of right translation
and is a subsemigroup, since $hp=kp=p$ implies
$(hk)p=h(kp)=hp=p$.  Minimality gives $E=S'$, and $p\in E$ proves
$p^2=p$.
\end{proof}

Set
\begin{equation}
\label{eq:map-closures}
 A_m=\overline{\conv\mathcal H_m}^{\,\mathrm{pw}},
 \qquad A=\bigcap_{m\geq1}A_m.
\end{equation}

\begin{proposition}[An anchor point]
\label{prop:anchor}
The set $A$ is a nonempty compact semigroup of nonexpansive selfmaps
of $K$ and contains an idempotent $P$.  There is a point $x\in K$ with
$P(x)=x$ such that the nonempty weakly compact set
\begin{equation}
\label{eq:anchor-tail-set}
 D_x=\bigcap_{m\geq1}
       \overline{\{h(x):h\in\mathcal H_m\}}^{\,\mathrm{weak}}
\end{equation}
satisfies
\begin{equation}
\label{eq:anchor-convex-hull}
 x\in\overline{\conv D_x}^{\,\norm{\cdot}}.
\end{equation}
\end{proposition}

\begin{proof}
The sets $A_m$ are nonempty, compact, and decreasing.  Every member
is nonexpansive by Lemma~\ref{lem:finite-tuples}.

First, $\conv\mathcal H_m$ is closed under composition.  Indeed, if
\[
 g=\sum_{r=1}^s\lambda_r(R_{n_r}\circ G_r)
 \in\conv\mathcal H_m,
 \qquad k\in\conv\mathcal H_m,
\]
then $k\in\mathcal G$ and
\[
 g\circ k
 =\sum_{r=1}^s\lambda_r\bigl(R_{n_r}\circ(G_r\circ k)\bigr)
 \in\conv\mathcal H_m.
\]
This identity distributes precomposition over the outer convex
combination.

To pass to $A_m$, fix $h,k\in A_m$, arguments $a_1,\ldots,a_q\in K$,
and $\delta>0$.  Lemma~\ref{lem:finite-tuples} gives
$k_\delta,h_\delta\in\conv\mathcal H_m$ with
\[
 \max_i\norm{k_\delta(a_i)-k(a_i)}<\delta,
 \qquad
 \max_i\norm{h_\delta(k(a_i))-h(k(a_i))}<\delta.
\]
The second application uses the fixed arguments $k(a_i)$.
Nonexpansiveness of $h_\delta$ then gives
\begin{align*}
 &\norm{h_\delta(k_\delta(a_i))-h(k(a_i))}\\
 &\quad\leq
   \norm{h_\delta(k_\delta(a_i))-h_\delta(k(a_i))}
   +\norm{h_\delta(k(a_i))-h(k(a_i))}\\
 &\quad\leq\norm{k_\delta(a_i)-k(a_i)}
   +\norm{h_\delta(k(a_i))-h(k(a_i))}
   <2\delta.
\end{align*}
The composition $h_\delta\circ k_\delta$ belongs to
$\conv\mathcal H_m$.  Since the finite arguments and $\delta$ were
arbitrary, these simultaneous norm approximations imply
$h\circ k\in A_m$.

Thus $A$ is a nonempty compact semigroup, by the nested intersection
property.  For fixed $k\in A$, the right translation
$h\mapsto h\circ k$ is continuous in the pointwise weak topology:
its value at an argument $a$ is the evaluation of $h$ at the fixed
argument $k(a)$.  Lemma~\ref{lem:idempotent} supplies $P\in A$ with
$P\circ P=P$.  Choose any $a\in K$ and set $x=P(a)$; then $P(x)=x$.

Write
\[
 B_m=\overline{\{h(x):h\in\mathcal H_m\}}^{\,\mathrm{weak}}.
\]
These are nonempty nested weakly compact subsets of $K$, so
$D_x=\bigcap_m B_m$ is nonempty and weakly compact.  Since
$P\in A_m$, Lemma~\ref{lem:finite-tuples} at the single argument $x$
shows that
\[
 x=P(x)\in
 \overline{\conv\{h(x):h\in\mathcal H_m\}}^{\,\norm{\cdot}}
 \subseteq\overline{\conv B_m}^{\,\norm{\cdot}}.
\]
Consequently, for each $\ell\in X^*$,
\[
 \max_{z\in B_m}\ell(z)\geq\ell(x).
\]
The maximum exists because $B_m$ is weakly compact.  The sets
$B_m\cap\{z:\ell(z)\geq\ell(x)\}$ are nonempty nested compact sets;
their intersection gives a point of $D_x$ at which
$\ell\geq\ell(x)$.  Hence
\[
 \sup_{z\in D_x}\ell(z)\geq\ell(x)\qquad(\ell\in X^*).
\]
If \eqref{eq:anchor-convex-hull} failed, separation from the
norm-closed convex hull of $D_x$ would give a functional violating
this inequality.  This proves the proposition.
\end{proof}

\subsection{Successive choices with prescribed thresholds}

Fix the point $x$ supplied by Proposition~\ref{prop:anchor}.  Its
convex-hull property permits each new map to be selected at a scale
that depends on all preceding choices.

\begin{lemma}[Adaptive finite selection]
\label{lem:adaptive-selection}
Let $\tau>0$.  There is a successive selection procedure with the
following property.  Before the $i$th map is chosen, prescribe any
finite integer $m_i\geq1$, allowing this prescription to depend on the
entire previous history and on any auxiliary data already fixed at
that stage.  The procedure chooses $G_i\in\mathcal H_{m_i}$ and
terminates after finitely many steps, say $M$, with nonnegative
coefficients $\beta_1,\ldots,\beta_M$ such that
\begin{equation}
\label{eq:adaptive-return}
 \sum_{i=1}^M\beta_i=1,
 \qquad
 \norm*{\sum_{i=1}^M\beta_iG_i(x)-x}<\tau.
\end{equation}
The coefficients are assigned on termination.  Removing zero
coefficients gives a nonempty list with positive coefficients,
in the original order of selection.
\end{lemma}

\begin{proof}
Choose $d_1,\ldots,d_q\in D_x$ and nonnegative coefficients
$\lambda_1,\ldots,\lambda_q$ summing to $1$ such that
\[
 \norm*{\sum_{r=1}^q\lambda_rd_r-x}<\frac\tau2.
\]
Fix a metric $\rho$ inducing the weak topology on $K$.  Cycle through
the targets $d_1,\ldots,d_q$: at step $i$, let $r(i)$ be the member of
$\{1,\ldots,q\}$ congruent to $i$ modulo $q$.  Once the threshold
$m_i$ has been prescribed, the membership
\[
 d_{r(i)}\in
 \overline{\{h(x):h\in\mathcal H_{m_i}\}}^{\,\mathrm{weak}}
\]
allows a choice $G_i\in\mathcal H_{m_i}$ with
\[
 \rho\bigl(G_i(x),d_{r(i)}\bigr)<\frac1i.
\]
Any auxiliary choices associated with this map may now be made before
the next threshold is prescribed.  Stop as soon as the convex hull
of the selected images comes within distance $\tau$ of $x$.

Suppose that the procedure never stops, and let
\[
 C_0=\overline{\conv\{G_i(x):i\geq1\}}^{\,\norm{\cdot}}.
\]
For each $r$, the subsequence assigned to target $d_r$ converges
weakly to $d_r$.  The set $C_0$ is weakly closed by separation, so
it contains every $d_r$ and their convex combination
$\sum_r\lambda_rd_r$.  By the definition of $C_0$, a finite convex
combination of the selected images is within $\tau/2$ of this last
point.  It is therefore within $\tau$ of $x$.  All its indices occur
in some finite prefix of the selection, contradicting the stopping
rule.  The procedure thus terminates, and a convex combination at
the stopping step supplies \eqref{eq:adaptive-return}.  Deleting its
zero coefficients leaves the stated positive list.
\end{proof}

By \eqref{eq:tail-defect}, the threshold at each selection may impose
any positive upper bound on the defect of that individual image.
The convex coefficients in Lemma~\ref{lem:adaptive-selection} provide
the return to $x$ in norm.

Finally, fix once and for all the sequence
\begin{equation}
\label{eq:terminal-sequence}
 b_j=R_j(x)\qquad(j\geq1).
\end{equation}
Lemma~\ref{lem:resolvents} and \eqref{eq:map-displacement} give
\begin{equation}
\label{eq:terminal-displacement}
 \norm{Tb_j-b_j}\leq\frac1j,
 \qquad
 \norm{G(b_j)-b_j}\leq\frac{c_G}{j}
 \quad(G\in\mathcal G).
\end{equation}
Thus every fixed finite family of maps in $\mathcal G$ acts
arbitrarily close to the identity at a common sufficiently late
$b_j$.  This sequence will be used throughout the remaining
construction.
